\documentclass[12pt]{extarticle}
\usepackage[utf8]{inputenc}
\usepackage[spanish]{babel}
\usepackage{cite}
\usepackage{fancyhdr}
\usepackage{graphicx}
\usepackage[a4paper,margin=1in,headheight=57 pt,headsep=12pt,heightrounded]{geometry}
\usepackage{titling}
\usepackage{lmodern}
\usepackage{parskip}
\usepackage{xcolor}

\definecolor{gris}{RGB}{220,220,220}

\date{\today}

\fancyhf{}
\pagestyle{fancy}
\rhead{ID: AEC-INF-AMB-001\\Versión: 1.1 \\Fecha: \thedate \\ Página: \thepage }
\lhead{\begin{picture}(0,0) \put(0,0){\includegraphics[width=40mm]{AEC.PNG}} \end{picture}}
\renewcommand{\headrulewidth}{0.5pt}
\chead{REPORTE CUBESAT \\ SOCIAL }





\begin{document}

\begin{titlepage}
  \vspace*{4cm}
  {\fontsize{18}{30}\selectfont\bfseries CUBESAT SOCIAL}
  \hfill
  % Logo
  \includegraphics[height=3cm]{AEC.PNG} \\
  % Línea gris
  {\color{gris}\hrule}
  \Large{\itshape Agencia Espacial De Colombia}
  \vfill
  \centering
  {\Huge \large Dylan Nicolás Sáenz Chavarro$^{1*}$}
  \vfill
  {\large $^{1}$dsaenzc@eafit.edu.co - practicante@agenciaespacialdecolombia.org  \hfill \thedate }
\end{titlepage}



\tableofcontents
\newpage

\begin{abstract}
  El presente documento informa sobre los análisis realizados en cuanto a las aplicaciones técnicas y análisis de misión necesarios en la implementación de un satélite de pequeño formato CubeSat en Colombia buscando solucionar problemáticas sociales del país.  
\end{abstract}



\section{Proyectos similares encontrados}
 Se realizó una búsqueda por medio de la base de datos $nanosats.eu$
de proyectos satélite de tipo CubeSat en todo el mundo que implementarán el desarrollo tecnológico de la misión empleada en solucionar problemáticas ambientales en ciudades, de esta búsqueda se encontraron siete proyectos satélite de interés para la agencia.
\begin{itemize}
    \item Conasat (A,B) - Brasil
    \item NITESat - Estados Unidos
    \item We Wish - Japón
    \item GEMS 2 - Reino unido
    \item MeznSat - Emiratos Arabes Unidos
    \item MOCI - Estados Unidos
\end{itemize}

\section{Conasat}

Este proyecto tiene estimado el lanzamiento en el año 2020 y buscará realizar toma de datos ambientales como volumen de lluvia, temperatura, humedad, contaminación del aire, corrientes oceánicas y riesgos ambientales.los datos son recopilados y transmitidos al satélite por plataformas remotas en tierra, y retransmitidos al centro de la mando.

\subsection{Características}
Distribución de tamaño 
\begin{itemize}
\item 2U ocupan todos los subsistemas
\item 2U redundancia de frío 
\item 4U superficie para paneles solares    
\end{itemize}

Sensores
\begin{itemize}
\item Sensores análogos solares
\item Magnetómetros 3 ejes 
\item Giróscopos de estructura vibrante
\end{itemize}

Actuadores
\begin{itemize}
\item Ruedas de reacción
\item ADSC Magnetorquer
\end{itemize}

\section{NITESat}

Este satélite adquiere imágenes nocturnas de alta calidad del Medio Oeste desde la órbita para medir la cantidad y calidad de la contaminación lumínica en un área de 1000 km x 1000 km centrada en Chicago. El objetivo secundario de la misión demostrará un refinamiento y monitoreo útiles de elementos de dos líneas (TLE) a través de un subsistema de repetidor de radio económico, de bajo volumen y baja potencia para satélites pequeños, trabajando en conjunto con una o más estaciones terrestres.

\subsection{Características}
Payload
\begin{itemize}
    \item Camara RGB pco.edge 3.1 2048 x 1536 pixeles 
\end{itemize}

Apoyo en tierra

\begin{itemize}
    \item GPS
    \item Magnetómetro
    \item Cámara  Blackfly GigE Aptina MT9P006
\end{itemize}

\section{WE WISH}
Construido por miembros del Meisei Amateur Radio Club en Japón. El satélite llevó una cámara infrarroja para estudios ambientales atmosféricos empleando imágenes de 320 por 256 píxeles de la superficie de la Tierra que fueron transmitidas en aproximadamente 110 segundos utilizando SSTV. 

\subsection{Características}
Sensores
\begin{itemize}
    \item Cámara infrarroja 320x256 píxeles
    \item GPS Radiosonde
\end{itemize}
\subsection{GPS Radiosonda}
Es un instrumento de sondeo en la parte superior del aire para medir varios tipos de datos meteorológicos; velocidad del viento, dirección del viento, presión, temperatura y humedad. La velocidad del viento, la dirección y la presión del viento se calculan a partir de la velocidad de desplazamiento y la altitud obtenida por las técnicas de posicionamiento GPS; y datos de temperatura y humedad medidos desde un termistor sintonizado y un sensor de humedad de capacitancia electrostática respectivamente.
\section{GEMS}
 Los satélites se lanzarán desde la Estación Espacial Internacional, lo que permitirá la demostración en órbita de propuestas técnicas y comerciales. Orbital Micro Systems se anunció como la primera misión de demostración en órbita (IOD), poniendo en el espacio una nueva tecnología de observación y pronóstico del tiempo en miniatura.
 
 \section{MeznSat}
 Lleva un espectrómetro infrarrojo de onda corta (SWIR) como carga útil principal, con el objetivo de sensar la cantidad de dióxido de carbono (CO2) y metano (CH4) y Agua (H2O) en la atmósfera haciendo observaciones dentro del espectro de onda corta 900-1700 nm, sobre EAU. Esta misión sigue misiones anteriores exitosas como CANX-2 y SathyabhamaSat que llevaban la misma carga útil (espectrómetro Argus 2000).
 
 \subsection{Características}
 Payload
 \begin{itemize}
     \item espectrómetro Argus 2000
     \item RGB cámara 1280x1024 píxeles 
 \end{itemize}
 Sensores
 \begin{itemize}
     \item ADCS Magnetómetro
     \item Placa ADCS de 3 ejes con sensor solar, sensor de tierra,
Magnetorquer, Giroscopio, CubeCoil
\item ADCS Rastreador de estrellas
 \end{itemize}
 Actuadores
 \begin{itemize}
     \item Rueda de reacción de 3 ejes
 \end{itemize}
 
 \section{MOCI}
 Realiza la estructuras desde el movimiento (SfM) en  órbita terrestre baja (LEO) y genera nubes de puntos 3D de estructuras a gran escala en la superficie de la Tierra. Es la primera vez que un CubeSat se ha especializado en la construcción de modelos 3D utilizando la estructura del movimiento. Busca que los investigadores pueden usar estos datos para hacer mejores predicciones sobre los modelos climáticos, los impactos ambientales, medir los recursos naturales y distinguir los cambios geográficos a lo largo del tiempo. La misión secundaria del MOCI es recopilar datos de color del océano de la órbita terrestre baja y rastrear la salud de la costa de Georgia. Estos datos ayudarán a los investigadores a evaluar mejor cómo el cambio climático está afectando la costa de Georgia. También es posible utilizar las capacidades SfM de MOCI para detectar el cambio relativo en el nivel del agua de varias costas y lagos.
 \subsection{Características}
 Payload
 \begin{itemize}
     \item Ocean Color Imager (6 visible, 2 NIR)
 \end{itemize}
 \section{¿GOES?}
  El programa fue diseñado para operar en órbita geoestacionaria esto implica que mantiene un punto fijo respecto a la tierra, a 35.790 km, Se encuentra continuamente a la vista de EE. UU. continental, de los océanos Pacífico y Atlántico, América Central y Sudamérica.
Permite establecer un monitoreo en nubes, temperatura superficial y vapor de agua, y sondeando los perfiles verticales de estructuras térmicas y de vapor. Así sigue la evolución de fenómenos de la atmósfera, asegurando la cobertura en tiempo real para seguir eventos de corta vida, especialmente tormentas locales, ciclón tropical, que directamente afectan la seguridad pública, protección de propiedades, y últimamente, salud y desarrollo económico. La importancia de esta capacidad ha sido recientemente ejemplificada durante los huracanes Hugo (1989) y Andrew (1992).


En este programa se cuenta con varios satélites  pero ninguno de ellos bajo el formato Cubesat un alternativo en la construcción de satélites de pequeño tamaño. Sin embargo, este cuenta con varias aplicaciones que pueden ser empleadas en cualquier tipo de monitoreo satelital como:
\begin{itemize}
    \item Demostración tecnológica 
    \item Seguimiento remoto de la tierra
    \item Efectos del entorno espacial en sistemas biológicos
    \item Efectos de radiación en tecnologías espaciales
    \item Astrofísica 
    \item Detección de terremotos (midiendo frecuencias extremadamente bajas de señales magnéticas)
    \item Servicios de telecomunicaciones, almacenamiento de datos, distribuciones de servicios, control y procesamiento de datos.
\end{itemize}

Este tipo de programas resulta muy importante no solo para los países desarrolladores sino para terceros  un claro ejemplo es el permiso que solicito Japón en el año 2003 del uso de las imágenes por uno de los satélites GOES para monitorear las zonas costeras por riesgo de catástrofes marítimas.

\section{Como transmitir datos}

 lo más común en las comunicaciones satélite es emplear el uso de ondas electromagnéticas que son capaces de viajar sin medio de propagación entre estas ondas se encuentra La radio y el espectro visible. los satélites de pequeño formato frecuentemente recurren al uso de comunicación por  radiofrecuencia empleando antenas UHF y VHF en las naves espaciales, o también por medio de banda $S$ que abarca frecuencias de 2025 a 2120 Mega Hertz. 
 
Estos datos se pueden transmitir al centro de mando desde la estación de radio terrestre dependiendo de la frecuencia de actualización de cada satélite, es frecuente enviar datos vía
 usb o dvd si se tiene una latencia de 6-48 H, internet para latencias en un rango de minutos,  de otro modo datos con mucho volumen o  varios requerimientos de control requieren un servidor web en donde almacenar a tiempo real para esto es necesario un lugar con alto ancho de banda.
 
 \subsection{Paquetes de datos}
Los datos pueden ser catalogados según el tipo de información científica que manejen, una clasificación reconocida y la más común es la que se emplea en NASA una clasificación por niveles según el tratamiento de los datos.
\begin{itemize}
    \item Nivel 0 : Son datos científicos de instrumentos de medida en resolución completa esto implica que son datos sin procesar por ejemplo lecturas de voltajes y contadores también se pueden encontrar ordenados temporalmente, sin errores de duplicado o perdida de transmisión. 
    \item Nivel 1A : Datos tratados de forma reversible, empaquetados y con datos auxiliares ej: datos resultantes a la  aplicación de ecuaciones. 
    
    \item Nivel 1B : transformación irreversible (resampleado,remapeado,calibrado) ej: Radianzas, intensidad de campo magnético,etc.
    \item Nivel 1C : Datos 1A- 1B que han sido remapeados en unas rejillas uniformes de espacio tiempo esto significa datos
    Radio-métricamente correctos inclusive también tienen correcciones adicionales como corrección de terreno.
    \item Nivel 2 : parámetros geofísicos, tomados del nivel 1, con un instrumento de ubicación posicionamiento y muestreo.
    \item Nivel 3: Parámetros geofísicos mapeados y ubicados en rejillas uniformes de espacio tiempo.  
\end{itemize}

\section{Autonomía del satélite}
\begin{itemize}
    \item Lazo cerrado: no requiere entrada de estación en tierra ej: control térmico del satélite
    \item Lazo cerrado abierto: establecer una configuración segura luego de una falla y esperar por la respuesta de tierra
    \item  Lazo cerrado: restablecer configuración nominal de la misión seguido de una falla y toma resumen de ejecución de procesos nominales.
\end{itemize}




\nocite{GOCI}
\nocite{MOCI}
\nocite{EAU}
 
 
 
 
\bibliographystyle{plain}
\bibliography{M335}

\end{document}
